\begin{table}[ht!]
\centering
\caption{How exposed to generative AI the occupations of Swedish women and men are, by age, 2024}
\label{tab:occ_mix_by_sex}
\footnotesize
\begin{tabular}{lcccc}
\toprule
 & \multicolumn{2}{c}{Mean exposure percentile} & \multicolumn{2}{c}{In the top quartile (\%)} \\
\cmidrule(lr){2-3}\cmidrule(lr){4-5}
Age & Women & Men & Women & Men \\
\midrule
16--24 & 39.0 & 34.8 & 10.3 & 9.3 \\
25--29 & 52.8 & 45.2 & 28.3 & 22.7 \\
30--34 & 56.7 & 49.8 & 33.6 & 28.4 \\
40--44 & 57.4 & 52.2 & 34.2 & 30.3 \\
50--54 & 57.1 & 52.2 & 33.8 & 29.3 \\
\midrule
All 16--64 & 54.3 & 48.6 & 30.3 & 25.9 \\
\bottomrule
\end{tabular}
\begin{minipage}{0.92\textwidth}\footnotesize\vspace{4pt}
Employees aged 16 to 64 in Sweden in 2024, by four-digit SSYK 2012 occupation, age and sex, from Statistics Sweden's occupational register (table YREG54BAS); 4,802,448 of 4,838,418 employees hold an occupation the DAIOE generative-AI index prices. The exposure percentile is that index, the same one the employment design uses, and the top quartile is its own. Each column is the employment-weighted mean over the group's occupational distribution. This is the NATIONAL distribution: part of the gap it shows is women and men working at different employers, which the within-employer design absorbs, and the table cannot say how much of the rest operates inside a firm. It states a difference in the occupations the two sexes hold; it does not decompose the estimated differential.
\end{minipage}
\end{table}
